\documentstyle[%


righttag%


,11pt%


,amssymb%


,russian%


,izvru%


]{amsart}





\newcommand{\tl}[3]{\medskip\noindent{\bf#1}\ #2 \dotfill #3 \par} 


\pagestyle{empty}


\begin{document}


%\tableofcontents


\par


\vskip 2cm


\par


\bigskip





Данный тематический номер посвящается светлой памяти заслуженного деятеля


науки Российской Федерации и Республики Татарстан, доктора физико-математических


наук, первой в истории Казанского университета женщины--профессора


математики, долгие годы возглавлявшей кафедру дифференциальных уравнений


Любови Ивановны Чибриковой.





\bigskip


\bigskip


\par


\centerline{\Large СОДЕРЖАНИЕ}


\bigskip


\bigskip


%%%%%%%%%% Use \newline %%%%%%%%%%%











\tl{Авхадиев Ф.Г.}


{Новые изопериметрические неравенства


для моментов областей и\newline жесткости кручения}{3}


\tl{Габдулхаев Б.Г.}


{Проекционные методы решения сингулярных интегральных уравнений}{12}


\tl{Гарифьянов Ф.Н., Салехова И.Г.} 


{Квазипериодическая краевая задача Римана в случае\newline  


переменного коэффициента}{25}


\tl{Елизаров А.М., Лапин А.В.}


{Применение вариационных методов в обратных


краевых\newline задачах для аналитических функций}{30}


\tl{Жегалов В.И.}


{О случаях разрешимости гиперболических уравнений


в квадратурах}{47}


\tl{Обносов Ю.В.}


{Решение задачи $\Bbb R$-линейного сопряжения в случае


гиперболической линии\newline разделения разнородных фаз}{53}


\tl{Плещинский Н.Б., Тумаков Д.Н.}


{Граничные задачи для уравнения Гельмгольца


в\newline квадранте и в полуплоскости,


составленной из двух квадрантов}{63}


\tl{Салехов Л.Г., Салехова Л.Л.}


{О некоторых классах уравнений в сверточной алгебре $D'_+$.\;.}{75}


\tl{Сильвестров В.В.} 


{Метод римановых поверхностей в задаче о межфазных трещинах и\newline 


включениях при наличии сосредоточенных сил}{78}

















\vfill


\noindent\raisebox{2pt}{\copyright} Казанский госудаpственный унивеpситет


\end{document}


